% 数学ノート集・共通テンプレート v2（積分論2026年度第1部の版面を参考）
% A5判。章内で定義・定理・例・注意等の番号を共通化。
% コンパイル： platex -interaction=nonstopmode math_note_template_v2.tex (2回)
%              dvipdfmx math_note_template_v2.dvi
\documentclass[a5paper,10pt,onecolumn]{jsarticle}
\usepackage[dvipdfmx,top=22mm,bottom=20mm,inner=19mm,outer=18mm,headheight=13pt,headsep=6mm]{geometry}
\usepackage{amsmath,amssymb,amsthm,mathtools}
\usepackage{enumitem,fancyhdr,titlesec}
\usepackage[dvipdfmx,bookmarks=true,colorlinks=false]{hyperref}
\usepackage{pxjahyper}

% ----- 差し替える文書情報 -----
\newcommand{\NoteTitle}{条件付き期待値とマルチンゲール}
\newcommand{\NoteSeries}{数学ノート集}

% ----- 文字・組版 -----
% 参考PDFの質素な明朝＋Computer Modernの印象を引き継ぎながら，
% A5判で窮屈にならないよう本文の行間を少し確保する．
\linespread{1.13}
\setlength{\parindent}{1zw}
\setlength{\parskip}{0pt}
\setlength{\abovedisplayskip}{6pt plus 2pt minus 2pt}
\setlength{\belowdisplayskip}{6pt plus 2pt minus 2pt}
\setlength{\abovedisplayshortskip}{3pt plus 2pt minus 1pt}
\setlength{\belowdisplayshortskip}{4pt plus 2pt minus 1pt}
\setlist[enumerate]{topsep=3pt,itemsep=1pt,parsep=0pt,leftmargin=2.4zw}
\titleformat{\section}{\normalfont\normalsize\bfseries}{\thesection}{0.8zw}{}
\titlespacing*{\section}{0pt}{2.0ex plus .4ex}{.9ex}
\titleformat{\subsection}{\normalfont\normalsize\bfseries}{\thesubsection}{0.7zw}{}
\titlespacing*{\subsection}{0pt}{1.5ex plus .3ex}{.6ex}
\raggedbottom

% ----- 柱とページ番号：ページ上部、罫線なし -----
\fancypagestyle{mathnote}{%
  \fancyhf{}%
  \fancyhead[L]{\footnotesize \NoteSeries\ （\NoteTitle）}%
  \fancyhead[R]{\footnotesize \thepage}%
  \renewcommand{\headrulewidth}{0pt}%
  \renewcommand{\footrulewidth}{0pt}%
}
\pagestyle{mathnote}

% ----- 節内通し番号の定義・定理群 -----
\newtheoremstyle{mathnotetheorem}
 {4.5pt}{4.5pt}{\normalfont}{}%
 {\bfseries}{}{0.65zw}%
 {\thmname{#1}\thmnumber{ #2}\thmnote{（#3）}}
\theoremstyle{mathnotetheorem}
\newtheorem{theorem}{定理}[section]
\newtheorem{definition}[theorem]{定義}
\newtheorem{lemma}[theorem]{補題}
\newtheorem{proposition}[theorem]{命題}
\newtheorem{corollary}[theorem]{系}
\newtheorem{example}[theorem]{例}
\newtheorem{remark}[theorem]{注意}
\renewcommand{\proofname}{証明}
\renewcommand{\qedsymbol}{$\square$}
\makeatletter
\renewenvironment{proof}[1][\proofname]{%
  \par\pushQED{\qed}\normalfont\topsep4\p@\@plus2\p@\relax
  \trivlist\item[\hskip\labelsep\bfseries #1]\ignorespaces
}{\popQED\endtrivlist\@endpefalse}
\makeatother

% ----- 使うノートで適宜拡張する記号 -----
\newcommand{\R}{\mathbb{R}}
\newcommand{\norm}[1]{\lVert #1\rVert}
\newcommand{\abs}[1]{\lvert #1\rvert}

\begin{document}
\section*{読了後にできること}
本稿を読むための前提は，条件付き期待値，$sigma$-代数と積分の定義．
本稿の具体的な到達課題は，有限分割に関する条件付き期待値を自分で計算する．

\section{確率空間と可積分確率変数}

$(\Omega,\mathcal F,P)$ を確率空間とし，
$\mathcal G\subset\mathcal F$ を部分 $\sigma$-加法族とする．
以下，確率変数は実数値とする．
$L^p(\mathcal F)$ は $P$ に関する同値類の空間を表し，
$E[X]=\int_\Omega X\,dP$ と書く．

\begin{definition}
$X\in L^1(\mathcal F)$ に対し，
$Y\in L^1(\mathcal G)$ が
\[
 \int_A Y\,dP=\int_A X\,dP\qquad(A\in\mathcal G)
 \tag{$*$}
\]
を満たすとき，
$Y$ を $X$ の $\mathcal G$ に関する\emph{条件付き期待値}といい，
$E[X\mid\mathcal G]$ と書く．等号は概同値類として理解する．
\end{definition}

\begin{theorem}[条件付き期待値の存在と一意性]\label{thm:ce}
任意の $X\in L^1(\mathcal F)$ に対し，
$E[X\mid\mathcal G]$ が存在し，ほとんど至る所で一意である．
\end{theorem}
\begin{proof}
$X\ge0$ の場合を考える．
$A\in\mathcal G$ に対して
$\nu(A)=\int_A X\,dP$ と置くと，
$\nu$ は $\mathcal G$ 上の有限正測度であり，
$\nu\ll P|_{\mathcal G}$ である．
Radon--Nikodymの定理により
$Y\in L^1(\mathcal G)$，$Y\ge0$ が存在して
$\nu(A)=\int_A Y\,dP$ を満たす．

一般の $X$ について $X=X^+-X^-$ と分解し，
それぞれに対応する $Y^+,Y^-$ を取る．
$Y=Y^+-Y^-\in L^1(\mathcal G)$ は $(*)$ を満たす．

$Y,Z$ が $(*)$ を満たすなら，
$W=Y-Z$ は $\mathcal G$ 可測で
$\int_AW\,dP=0$ が全ての $A\in\mathcal G$ で成り立つ．
$A_n=\{W>1/n\}$ を取ると
$0=\int_{A_n}W\,dP\ge P(A_n)/n$ だから $P(A_n)=0$ である．
$\{W>0\}=\bigcup_nA_n$ なので $P(W>0)=0$ となる．
同様に $P(W<0)=0$ である．ゆえに $Y=Z$ 概至る所である．
\end{proof}

\begin{proposition}[基本性質]\label{prop:ce-basic}
$X,Y\in L^1(\mathcal F)$，$a,b\in\mathbb R$ とする．
\begin{enumerate}
\item $E[aX+bY\mid\mathcal G]
 =aE[X\mid\mathcal G]+bE[Y\mid\mathcal G]$．
\item $X\ge0$ なら $E[X\mid\mathcal G]\ge0$．
特に $X\le Y$ なら
$E[X\mid\mathcal G]\le E[Y\mid\mathcal G]$．
\item $X$ が $\mathcal G$ 可測なら
$E[X\mid\mathcal G]=X$．
\item $E[E[X\mid\mathcal G]]=E[X]$．
\end{enumerate}
いずれも等式・不等式はほとんど至る所の意味である．
\end{proposition}
\begin{proof}
(1) は積分の線形性と定理\ref{thm:ce}の一意性から従う．
(2) は存在証明の正値性と (1) から従う．
(3) は $X$ 自身が定義の二条件を満たすからである．
(4) は $(*)$ に $A=\Omega$ を代入すればよい．
\end{proof}

\begin{lemma}[有界な因子の取り出し]\label{lem:pullout}
$X\in L^1(\mathcal F)$，
$Z$ を有界な $\mathcal G$ 可測確率変数とする．
このとき
\[
 E[ZX\mid\mathcal G]=Z\,E[X\mid\mathcal G].
\]
\end{lemma}
\begin{proof}
$Z=1_B$（$B\in\mathcal G$）の場合，
$A\in\mathcal G$ に対して
\[
 \int_A1_B E[X\mid\mathcal G]\,dP
 =\int_{A\cap B}X\,dP
 =\int_A1_BX\,dP
\]
なので一意性から主張が成り立つ．
単純 $\mathcal G$ 可測 $Z$ については線形性から従う．
一般の有界 $Z$ については
$\norm{Z_n-Z}_\infty\to0$ となる
単純 $\mathcal G$ 可測関数列 $(Z_n)$ を取る．
\[
 E[|Z_nX-ZX|]
 \le\norm{Z_n-Z}_\infty E|X|\to0
\]
であり，条件付き期待値の定義から
\[
 E[|E[W\mid\mathcal G]|]
 \le E|W|
 \tag{$**$}
\]
が成り立つ．実際，
$-E[|W|\mid\mathcal G]\le E[W\mid\mathcal G]
\le E[|W|\mid\mathcal G]$ に注意すればよい．
従って両辺の $L^1$ 極限を取って主張を得る．
\end{proof}

\begin{proposition}[塔の性質]\label{prop:tower}
$\mathcal H\subset\mathcal G\subset\mathcal F$ なら，
$X\in L^1(\mathcal F)$ に対して
\[
 E[E[X\mid\mathcal G]\mid\mathcal H]
   =E[X\mid\mathcal H]
\]
が成り立つ．
\end{proposition}
\begin{proof}
左辺は $\mathcal H$ 可測かつ可積分である．
任意の $A\in\mathcal H$ は $\mathcal G$ にも属するから
\[
 \int_A E[E[X\mid\mathcal G]\mid\mathcal H]\,dP
 =\int_AE[X\mid\mathcal G]\,dP
 =\int_AX\,dP.
\]
定理\ref{thm:ce}の一意性により結論を得る．
\end{proof}

\begin{proposition}[独立な情報]\label{prop:independent}
$X\in L^1(\mathcal F)$ が $\mathcal G$ と独立なら
$E[X\mid\mathcal G]=E[X]$ である．
\end{proposition}
\begin{proof}
$X$ と $1_A$ は独立である（$A\in\mathcal G$）から
$E[X1_A]=E[X]P(A)$ である．
定数 $E[X]$ は $\mathcal G$ 可測であり，
条件付き期待値の定義を満たす．
\end{proof}

\section{条件付きJensenの不等式}

\begin{lemma}[凸関数のアフィン下界]\label{lem:convex-support}
$\varphi:\mathbb R\to\mathbb R$ を有限値の凸関数とする．
各 $r\in\mathbb R$ に対し
$a_r\in\mathbb R$ が存在して
\[
 \varphi(x)\ge\varphi(r)+a_r(x-r)
 \qquad(x\in\mathbb R)
\]
を満たす．また有理数 $r$ のみを使って
\[
 \varphi(x)=\sup_{r\in\mathbb Q}
 \{\varphi(r)+a_r(x-r)\}
\]
と表せる．
\end{lemma}
\begin{proof}
凸性から割線の傾きは端点を右へ動かすと単調増加する．
従って
\[
 \sup_{u<r}\frac{\varphi(r)-\varphi(u)}{r-u}
 \le
 \inf_{v>r}\frac{\varphi(v)-\varphi(r)}{v-r}
\]
が成り立ち，
両辺は $r$ の近傍の二点を固定して評価すれば有限である．
この間から $a_r$ を取ると，
$r$ より左右いずれにある $x$ についても
上のアフィン下界が従う．
有限値の凸関数は実数直線上で連続である．
これは，$r$ の左右の固定した二点を用いて凸性による
局所的な傾きの上下界を得れば示せる．
任意の $x$ に有理数列 $r_n\to x$ を取ると，
$\varphi(r_n)+a_{r_n}(x-r_n)$ の
$x$ における値が $\varphi(x)$ に近づくことを示すには
$r_n$ を $x$ の近傍に限定すればよい．
凸関数の傾きはその近傍では一様有界なので
$a_{r_n}(x-r_n)\to0$ であり，
$\varphi(r_n)\to\varphi(x)$ となる．
従って有理数に関する上限は $\varphi(x)$ と一致する．
\end{proof}

\begin{theorem}[条件付きJensenの不等式]\label{thm:jensen}
$X\in L^1(\mathcal F)$ とし，
$\varphi:\mathbb R\to\mathbb R$ を有限値の凸関数とする．
$\varphi(X)\in L^1(\mathcal F)$ を仮定すると
\[
 \varphi(E[X\mid\mathcal G])
 \le E[\varphi(X)\mid\mathcal G]
 \qquad\text{概至る所}
\]
が成り立つ．
\end{theorem}
\begin{proof}
補題\ref{lem:convex-support}により，
各 $r\in\mathbb Q$ について
$\varphi(X)\ge\varphi(r)+a_r(X-r)$ である．
条件付き期待値の単調性と線形性から
\[
 E[\varphi(X)\mid\mathcal G]
 \ge\varphi(r)+a_r(E[X\mid\mathcal G]-r)
\]
が概至る所で成り立つ．
有理数全体は可算なので，同時に全ての $r\in\mathbb Q$
について成立する確率1の集合を取れる．
その集合上で $r\in\mathbb Q$ の上限を取れば
補題\ref{lem:convex-support}によって結論を得る．
\end{proof}

\begin{corollary}[$L^p$ 縮小性]\label{cor:lp}
$1\le p<\infty$，$X\in L^p(\mathcal F)$ とすると
\[
 E[|E[X\mid\mathcal G]|^p]\le E[|X|^p]
\]
が成立する．
\end{corollary}
\begin{proof}
$\varphi(t)=|t|^p$ は凸であるから
定理\ref{thm:jensen}と期待値の塔の性質を用いればよい．
\end{proof}

\section{フィルトレーションとマルチンゲール}

\begin{definition}
$\mathcal F_0\subset\mathcal F_1\subset\cdots\subset\mathcal F$
をフィルトレーションという．
確率過程 $(X_n)_{n\ge0}$ が適合であるとは，
$X_n$ が $\mathcal F_n$ 可測であることをいう．
適合な可積分過程 $(M_n)$ が\emph{マルチンゲール}であるとは，
\[
 E[M_{n+1}\mid\mathcal F_n]=M_n
\]
を満たすことをいう．
右辺を左辺が上回る場合を劣マルチンゲール，
逆の場合を優マルチンゲールという．
\end{definition}

\begin{example}[独立増分]
$\xi_1,\xi_2,\ldots$ を互いに独立な可積分確率変数とし，
$E\xi_k=0$ とする．
$M_n=\sum_{k=1}^n\xi_k$（$M_0=0$），
$\mathcal F_n=\sigma(\xi_1,\dots,\xi_n)$ と置けば
$(M_n)$ はマルチンゲールである．
\end{example}
\begin{proof}
$M_{n+1}=M_n+\xi_{n+1}$ であり，
$M_n$ は $\mathcal F_n$ 可測である．
独立性から
$E[\xi_{n+1}\mid\mathcal F_n]=E\xi_{n+1}=0$ なので
条件付き期待値の線形性によって
$E[M_{n+1}\mid\mathcal F_n]=M_n$ となる．
\end{proof}

\begin{proposition}\label{prop:martingale-time}
$(M_n)$ がマルチンゲールなら，
$m\ge n$ に対して
$E[M_m\mid\mathcal F_n]=M_n$ である．
また $(|M_n|^p)$ は $p\ge1$ で
$M_n\in L^p$ のとき劣マルチンゲールとなる．
\end{proposition}
\begin{proof}
最初の主張は塔の性質を繰り返して得る．
二番目は凸関数 $t\mapsto|t|^p$ に対する
条件付きJensenの不等式を用いれば
\[
 E[|M_{n+1}|^p\mid\mathcal F_n]
 \ge|E[M_{n+1}\mid\mathcal F_n]|^p=|M_n|^p
\]
である．
\end{proof}

\begin{definition}
$\{T\le n\}\in\mathcal F_n$ が
すべての $n$ で成立する
$\{0,1,\dots\}\cup\{\infty\}$ 値確率変数 $T$ を
停止時刻という．
$\mathcal F_T=\{A\in\mathcal F:
 A\cap\{T\le n\}\in\mathcal F_n\text{ for all }n\}$
と定める．
\end{definition}

\begin{theorem}[有界停止時刻の停止定理]\label{thm:optional}
$(M_n)$ をマルチンゲール，
$S\le T\le N$ を有界停止時刻とする．
このとき
\[
 E[M_T\mid\mathcal F_S]=M_S
\]
が成り立つ．
\end{theorem}
\begin{proof}
$M_T,M_S$ は有限個の可積分確率変数の和として
表されるため可積分である．また $M_S$ は
$\mathcal F_S$ 可測である．
任意の $A\in\mathcal F_S$ に対して
\[
 M_T-M_S=\sum_{k=1}^N
 (M_k-M_{k-1})1_{\{S<k\le T\}}
\]
と書ける．
$A\cap\{S<k\}\cap\{T\ge k\}$ は
$\mathcal F_{k-1}$ に属する．
実際，
$A\cap\{S\le k-1\}\in\mathcal F_{k-1}$，
$\{T\ge k\}=\{T\le k-1\}^{c}\in\mathcal F_{k-1}$
だからである．
従って
\begin{align*}
 E[(M_T-M_S)1_A]
 &=\sum_{k=1}^NE[(M_k-M_{k-1})
       1_{A\cap\{S<k\le T\}}]\\
 &=\sum_{k=1}^NE[
 E[M_k-M_{k-1}\mid\mathcal F_{k-1}]
       1_{A\cap\{S<k\le T\}}]=0.
\end{align*}
よって条件付き期待値の特徴づけから結論が従う．
\end{proof}

\section{Doobの最大不等式}

\begin{theorem}[弱型最大不等式]\label{thm:doob-weak}
$(X_k)_{k=0}^N$ を非負の劣マルチンゲールとし，
$\lambda>0$ とする．
\[
 A=\left\{\max_{0\le k\le N}X_k\ge\lambda\right\}
\]
と置くと
\[
 \lambda P(A)\le E[X_N1_A]\le E[X_N]
\]
が成り立つ．
\end{theorem}
\begin{proof}
$T=\min\{0\le k\le N:X_k\ge\lambda\}$ とし，
該当する時刻がなければ $T=N$ とする．
$A\cap\{T=k\}$（$0\le k\le N$）は
$\mathcal F_k$ に属する．
劣マルチンゲールの性質と塔の性質から
$E[X_N\mid\mathcal F_k]\ge X_k$ である．
従って
\begin{align*}
 E[X_N1_A]
 &=\sum_{k=0}^NE[X_N1_{A\cap\{T=k\}}]\\
 &\ge\sum_{k=0}^NE[X_k1_{A\cap\{T=k\}}]
 =E[X_T1_A]\ge\lambda P(A).
\end{align*}
最後の不等式は $X_N\ge0$ から従う．
\end{proof}

\begin{theorem}[Doobの $L^p$ 最大不等式]\label{thm:doob-lp}
$p>1$ とし，$(M_k)_{k=0}^N$ を
$L^p$ 可積分なマルチンゲールとする．
このとき
\[
 E\left[\max_{0\le k\le N}|M_k|^p\right]
 \le \left(\frac{p}{p-1}\right)^p E|M_N|^p.
\]
\end{theorem}
\begin{proof}
$Y_k=|M_k|$ は非負の劣マルチンゲールである．
$Y^*=\max_{0\le k\le N}Y_k$ と置く．
定理\ref{thm:doob-weak}により
$\lambda P(Y^*\ge\lambda)
 \le E[Y_N1_{\{Y^*\ge\lambda\}}]$ である．
$p\lambda^{p-2}$ を掛けて $\lambda>0$ で積分し，
Tonelliの定理を用いると
\[
 E[(Y^*)^p]
 \le \frac{p}{p-1}E[Y_N(Y^*)^{p-1}].
\]
右辺はH\"olderの不等式により
\[
 \frac{p}{p-1}
 (E[Y_N^p])^{1/p}
 (E[(Y^*)^p])^{(p-1)/p}
\]
以下である．
$N$ は有限だから $E[(Y^*)^p]<\infty$ であり，
これが零でない場合には両辺を適切な冪で割れば
結論を得る．零の場合は自明である．
\end{proof}

\section{上方横断とマルチンゲール収束}

\begin{definition}
実数 $a<b$ と列 $x_0,\dots,x_N$ に対し，
$x_k\le a$ となる時刻から
$x_l\ge b$ となる後の時刻までの往復を一回の
\emph{上方横断}と数える．
このような互いに重ならない完了した上方横断の
最大回数を $U_N(a,b)$ と書く．
確率過程に対しては標本路ごとに定義する．
\end{definition}

\begin{theorem}[Doobの上方横断不等式]\label{thm:upcrossing}
$(X_n)_{n=0}^N$ を劣マルチンゲールとする．
任意の $a<b$ に対して
\[
 (b-a)E[U_N(a,b)]
 \le E[(X_N-a)^+]-E[(X_0-a)^+]
\]
が成り立つ．
\end{theorem}
\begin{proof}
$Y_k=(X_k-a)^+$ とすると，
条件付きJensenの不等式と単調性により
$(Y_k)$ は非負の劣マルチンゲールである．
時刻 $k$ の直前までの $X_0,\dots,X_{k-1}$ により，
直前に $X\le a$ を観測してから
まだ $X\ge b$ を観測していないとき
$H_{k-1}=1$，そうでないとき $H_{k-1}=0$ と定める．
$H_{k-1}$ は $\mathcal F_{k-1}$ 可測である．
各完了上方横断に対応する増分の和は，
開始時の $Y=0$ と終了時の $Y\ge b-a$ から
$b-a$ 以上である．
最後に未完了の横断がある場合にも，
その区間の増分の和は $Y_N-0\ge0$ である．
従って標本路ごとに
\[
 (b-a)U_N(a,b)
 \le\sum_{k=1}^N H_{k-1}(Y_k-Y_{k-1}).
 \tag{$*$}
\]
ここで
\[
 E[(1-H_{k-1})(Y_k-Y_{k-1})]\ge0
\]
が劣マルチンゲール性から従う．
よって
\[
 E\sum_{k=1}^NH_{k-1}(Y_k-Y_{k-1})
 \le E\sum_{k=1}^N(Y_k-Y_{k-1})
 =E[Y_N]-E[Y_0].
\]
$(*)$ に期待値を取れば所望の不等式を得る．
\end{proof}

\begin{theorem}[マルチンゲール収束定理の基本形]
\label{thm:mart-conv}
$(X_n)_{n\ge0}$ を劣マルチンゲールとし，
$\sup_nE|X_n|<\infty$ を仮定する．
このとき $X_n$ はほとんど確実に有限値の
確率変数 $X_\infty\in L^1$ に収束する．
ただし，この仮定だけでは一般に $L^1$ 収束は従わない．
\end{theorem}
\begin{proof}
任意の有理数 $a<b$ に対して
定理\ref{thm:upcrossing}により
\[
 (b-a)E[U_N(a,b)]
 \le E[(X_N-a)^+]\le\sup_nE|X_n|+|a|.
\]
$U_N(a,b)$ は $N$ に関して単調増加なので，
単調収束定理によって
$E[U_\infty(a,b)]<\infty$ を得る．
従って $U_\infty(a,b)<\infty$ が概至る所で成立する．
有理数の対 $(a,b)$ は可算だから，
全ての有理数 $a<b$ についてこの性質が同時に成立する
確率1の集合を取れる．
その集合で
$\liminf_nX_n<\limsup_nX_n$ なら，
その間に有理数 $a<b$ を選ぶことによって
無限回の上方横断が生じて矛盾する．
よって $X_n$ は拡張実数値の極限 $X_\infty$ に収束する．
Fatouの補題より
\[
 E[\liminf_n|X_n|]
 \le\liminf_nE|X_n|
 \le\sup_nE|X_n|<\infty
\]
なので極限は有限値であり，かつ $E|X_\infty|<\infty$ である．
\end{proof}

\begin{example}[$L^1$ 収束しない非負マルチンゲール]
$\xi_1,\xi_2,\ldots$ を独立で
$P(\xi_k=0)=P(\xi_k=2)=1/2$ とし，
$M_n=\prod_{k=1}^n\xi_k$，$M_0=1$ と置く．
$(M_n)$ は非負マルチンゲールで $EM_n=1$ だが，
$M_n\to0$ はほとんど確実に成立する．
\end{example}
\begin{proof}
独立性と $E\xi_k=1$ から
$E[M_{n+1}\mid\mathcal F_n]=M_n$ となる．
$M_n$ が零でないのは最初の $n$ 個全てが $2$ の場合に限り，
その確率は $2^{-n}$ である．
この確率は $0$ に収束し，一度零になれば以後零なので，
$M_n\to0$ はほとんど確実に成立する．
もし $L^1$ 収束するなら期待値も収束して
$EM_n\to0$ となるはずだが，$EM_n=1$ なので矛盾する．
\end{proof}

\section{一様可積分性と $L^p$ 収束}

\begin{definition}
$\mathcal H\subset L^1(\mathcal F)$ が
\emph{一様可積分}であるとは，
\[
 \lim_{K\to\infty}
 \sup_{X\in\mathcal H}E[|X|1_{\{|X|>K\}}]=0
\]
を満たすことをいう．
\end{definition}

\begin{lemma}[一様可積分収束]\label{lem:ui-vitali}
$(X_n)$ が一様可積分で，
$X_n\to X$ が概至る所で成り立つなら，
$X\in L^1$ かつ $E|X_n-X|\to0$ となる．
\end{lemma}
\begin{proof}
一様可積分性から $\sup_nE|X_n|<\infty$ である．
実際，任意の固定した $K>0$ について
$E|X_n|\le K+E[|X_n|1_{\{|X_n|>K\}}]$ である．
Fatouの補題より $E|X|<\infty$ を得る．
任意の $K>0$ に対し，
切断関数
$\tau_K(t)=\max\{-K,\min\{t,K\}\}$ を用いれば
より簡潔に
\[
 |X_n-X|
 \le|\tau_K(X_n)-\tau_K(X)|
    +|X_n|1_{\{|X_n|>K\}}
    +|X|1_{\{|X|>K\}}
\]
と評価できる．
最初の項の期待値は，
$|\tau_K(X_n)-\tau_K(X)|\le2K$ と
概収束からLebesgueの収束定理によって $0$ に向かう．
一様可積分性により第二項は $K\to\infty$ で
$n$ によらず $0$ に向かう．
第三項も $X\in L^1$ なので $K\to\infty$ で $0$ に向かう．
従って $E|X_n-X|\to0$ を得る．
\end{proof}

\begin{theorem}[一様可積分マルチンゲール]\label{thm:ui-mart}
$(M_n)$ をマルチンゲールとする．
次は同値である．
\begin{enumerate}
\item $(M_n)$ は一様可積分である．
\item ある $M_\infty\in L^1$ が存在して
$M_n\to M_\infty$ が概至る所かつ $L^1$ で成立する．
\item ある $Z\in L^1$ が存在して
$M_n=E[Z\mid\mathcal F_n]$ が各 $n$ について成り立つ．
\end{enumerate}
\end{theorem}
\begin{proof}
(1) なら一様可積分性から $\sup_nE|M_n|<\infty$ となる．
定理\ref{thm:mart-conv}により $M_n\to M_\infty$ は
概至る所で成立し，補題\ref{lem:ui-vitali}から $L^1$ でも成立する．

(2) とする．任意の $A\in\mathcal F_n$ と $m\ge n$ に対し
$E[M_m1_A]=E[M_n1_A]$ である．
$m\to\infty$ として $L^1$ 収束を用いれば
$E[M_\infty1_A]=E[M_n1_A]$ である．
従って $M_n=E[M_\infty\mid\mathcal F_n]$ となり (3) を得る．

(3) とする．$Z\in L^1$ に対して
$Z_K=Z1_{\{|Z|\le K\}}$，
$R_K=Z-Z_K$ と置く．
条件付き期待値の単調性により
\[
 |E[Z\mid\mathcal F_n]|
 \le K+E[|R_K|\mid\mathcal F_n].
\]
$A_n=\{|E[Z\mid\mathcal F_n]|>2K\}$ と置くと，
$A_n\in\mathcal F_n$ であり，
$A_n$ 上で $E[|R_K|\mid\mathcal F_n]>K$ となる．
従って
\[
 KP(A_n)\le E[|R_K|1_{A_n}]\le E|R_K|
\]
であり，
\begin{align*}
 E[|M_n|1_{A_n}]
 &\le KP(A_n)+E[|R_K|1_{A_n}]\\
 &\le 2E|R_K|.
\end{align*}
$K\to\infty$ とすると右辺は $0$ に収束するので，
$(M_n)$ は一様可積分である．
\end{proof}

\begin{corollary}[$L^p$ 有界マルチンゲールの収束]\label{cor:lp-mart}
$p>1$ とし $(M_n)$ をマルチンゲールとする．
$\sup_nE|M_n|^p<\infty$ なら，
ある $M_\infty\in L^p$ が存在して
$M_n\to M_\infty$ がほとんど確実にかつ
$L^p$ で成立する．
\end{corollary}
\begin{proof}
$C=\sup_nE|M_n|^p<\infty$ と置く．
$|M_n|1_{\{|M_n|>K\}}
 \le |M_n|^p/K^{p-1}$ より一様可積分である．
定理\ref{thm:ui-mart}によって概収束する
$M_\infty\in L^1$ が存在する．Fatouの補題から
$E|M_\infty|^p\le C$ である．
定理\ref{thm:doob-lp}を各有限区間 $[0,N]$ に適用し，
単調収束定理を用いれば
\[
 E\left[\sup_{n\ge0}|M_n|^p\right]
 \le\left(\frac p{p-1}\right)^p C<\infty.
\]
従って
$|M_n-M_\infty|^p
 \le2^p\sup_{k\ge0}|M_k|^p$
という可積分な優関数が存在する．
Lebesgueの収束定理から
$E|M_n-M_\infty|^p\to0$ を得る．
\end{proof}

\begin{remark}
条件付き期待値の構成はRadon--Nikodymの定理に，
マルチンゲール収束は上方横断の制御に依存する．
確率積分論に進む際には，これらの結果が
適合性と収束様式を保証するために用いられる．
\end{remark}

\section{到達課題の完全解答}
\noindent\textbf{条件付き期待値の明示的構成．}
\begin{proposition}
$([0,1],\mathcal B,\lambda)$ を確率空間とし，
$X(\omega)=\omega$ とする．
$\mathcal G$ を
$A=[0,1/2)$ と $B=[1/2,1]$ が生成する
$\sigma$-代数とすると
\[
E[X\mid\mathcal G]
=\frac14\mathbf1_A+\frac34\mathbf1_B
\quad\text{a.s.}
\]
である．
\end{proposition}
\begin{proof}
右辺 $Y$ は $\mathcal G$-可測で可積分である．
$A$ 上で
$\int_A X\,d\lambda=\int_0^{1/2}\omega\,d\omega=1/8$
であり，
$\int_A Y\,d\lambda=(1/4)(1/2)=1/8$．
$B$ 上では
$\int_BX\,d\lambda=3/8$，
$\int_BY\,d\lambda=(3/4)(1/2)=3/8$．
$\mathcal G=\{\varnothing,A,B,[0,1]\}$ なので
全ての $\mathcal G$-可測集合で積分が一致し，
条件付き期待値の定義を満たす．
\end{proof}
\noindent\textbf{到達確認．}
同じ $\mathcal G$ で $E[X^2\mid\mathcal G]$ を求めよ．
\par\noindent\textbf{解答．}
$A$ 上では
$2\int_0^{1/2}\omega^2\,d\omega=1/12$，
$B$ 上では
$2\int_{1/2}^1\omega^2\,d\omega=7/12$．
従って $(1/12)\mathbf1_A+(7/12)\mathbf1_B$ である．
\end{document}
